\documentclass[11pt,a4paper]{article}
\usepackage[utf8]{inputenc}
\usepackage[T1]{fontenc}
\usepackage{lmodern}
\usepackage[ngerman]{babel}
\usepackage{amsmath,amssymb,amsthm}
\usepackage{geometry}
\usepackage{xcolor}
\usepackage{booktabs}

\geometry{verbose,tmargin=2.5cm,bmargin=2.5cm,lmargin=2.5cm,rmargin=2.5cm}

\newtheorem{theorem}{Theorem}
\newtheorem{definition}{Definition}

\definecolor{brandblue}{RGB}{20, 50, 90}

\title{\textbf{\color{brandblue} Topologische Aharonov-Bohm-Analogie \& Epstein-Zeta-Solenoid im quaternional-oktonionischen $E_8$-Modell}\\[0.5em]
\large \textit{Systematische Dokumentation für das \#Energiedoku-Manuskript}}
\author{Computer Institut Bamberg}
\date{August 2026}

\begin{document}

\maketitle

\begin{abstract}
Dieses Dokument fasst die exakte mathematische Mechanik der Aharonov-Bohm-Analogie im Rahmen des \textbf{\#Energiedoku}-Modells zusammen. Die radiale $E_8$-Epstein-Zeta-Funktion $Z_{E_8}(s)$ weist bei $s=8$ eine meromorphe Polstelle auf, die im oktonionischen Raum $\mathbb{O}$ als ein 6-dimensionaler, vermiedener Defektunterraum $\Sigma_8$ der Kodimension 2 fungiert. Wir zeigen die analytische Faktorisierung von $Z_{E_8}(s)$, bestimmen die holomorphe Schutzzone $R_0 = 10$ zur nächsten Nullstelle und beschreiben die topologische Eindeutigkeit der induzierten Zeta-Holonomie.
\end{abstract}

\section{Analytische Faktorisierung der $E_8$-Epstein-Zeta-Funktion}

Die radiale Epstein-Zeta-Funktion des $E_8$-Gitters lässt sich über die Darstellung der Theta-Reihe $\Theta_{E_8}(\tau) = E_4(\tau)$ mit $r_8(2n) = 240 \sigma_3(n)$ direkt auf ein Produkt von Riemannschen Zeta-Funktionen zurückführen:
\begin{equation}
Z_{E_8}(s) = \sum_{v \in E_8 \setminus \{0\}} \frac{1}{|v|^s} = 240 \cdot 2^{-s/2} \, \zeta\left(\frac{s}{2}\right) \zeta\left(\frac{s-6}{2}\right)
\end{equation}

Aus dieser Darstellung folgen unmittelbar die beiden meromorphen Polstellen:
\begin{itemize}
    \item \textbf{Pol 1:} bei $\frac{s}{2} = 1 \implies s = 2$ (kompensiert durch triviale Nullstelle des zweiten Faktors).
    \item \textbf{Solenoid-Pol (Hauptdefekt):} bei $\frac{s-6}{2} = 1 \implies \mathbf{s = 8}$.
\end{itemize}

\section{Topologische Geometrie \& Oktonionische Projektion}

Unter der Einbettung der Spektralebene in den 8-dimensionalen oktonionischen Raum $\mathbb{O} \cong \mathbb{R}^8$ bestimmt die Polstelle $s=8$ die topologischen Rahmendaten des Eich-Systems:

\begin{definition}[Defektunterraum $\Sigma_8$ und Komplement $M_8$]
Der Wert $s=8$ definiert einen reell 6-dimensionalen Defektunterraum $\Sigma_8 \subset \mathbb{O}$ der Kodimension 2. Der komplementäre reguläre Raum 
\begin{equation}
M_8 = \mathbb{O} \setminus \Sigma_8
\end{equation}
besitzt für transversale Schnittflächen die fundamentale Gruppe $\pi_1(M_8) \cong \mathbb{Z}$, was den topologischen Träger für nicht-triviale Phasen-Wicklungen bildet.
\end{definition}

Die logarithmische Ableitung $\omega_Z = \frac{Z'_{E_8}(s)}{Z_{E_8}(s)} ds$ definiert eine Differentialform mit verschwindender Feldstärke im regulären Raum ($d\omega_Z = 0$), welche für geschlossene Umläufe $\gamma \subset M_8$ eine topologisch geschützte Holonomie induziert:
\begin{equation}
\operatorname{Hol}_\gamma = \exp\left( -\alpha \oint_\gamma \omega_Z \right) = \exp(2\pi i \alpha k), \quad k \in \mathbb{Z}
\end{equation}

\section{Analytische Bestimmung der Schutzzone $R_0 = 10$}

Die Nullstellenmenge $\mathcal{Z} = \mathrm{Zeros}(Z_{E_8})$ besteht aus zwei Klassen von Singularitäten:
\begin{enumerate}
    \item \textbf{Kritische Linien (Nicht-triviale Nullstellen):} 
    \begin{itemize}
        \item Aus $\zeta(s/2)$: Nullstellen bei $\rho_n^{(1)} = 1 + 2i\gamma_n$ ($\gamma_1 \approx 14{,}1347$).
        \item Aus $\zeta(\frac{s-6}{2})$: Nullstellen bei $\rho_n^{(2)} = 7 + 2i\gamma_n$.
    \end{itemize}
    Der euklidische Abstand der nächstgelegenen nicht-trivialen Nullstelle $\rho_1^{(2)}$ zum Solenoidpol $s=8$ beträgt:
    \begin{equation}
    |8 - \rho_1^{(2)}| = \sqrt{(8-7)^2 + (2 \cdot 14{,}1347)^2} = \sqrt{1 + 799{,}16} \approx 28{,}29
    \end{equation}

    \item \textbf{Reelle Achse (Triviale Nullstellen):}
    Die trivialen Nullstellen von $\zeta(z)$ bei $z \in \{-2, -4, -6, \dots\}$ erzeugen Nullstellen bei $s \in \{-4, -8, \dots\}$ sowie $s \in \{-2, -6, \dots\}$. 
    Die absolut am nächsten gelegene Nullstelle liegt bei $s = -2$.
\end{enumerate}

\begin{theorem}[Holomorpher Schutzzonen-Radius]
Der minimale euklidische Abstand $R_0$ in der Spektralebene vom Solenoidpol $s=8$ zur Nullstellenmenge $\mathcal{Z}$ wird durch die triviale Nullstelle bei $s=-2$ exakt festgelegt:
\begin{equation}
R_0 = \min_{\rho \in \mathcal{Z}} |8 - \rho| = |8 - (-2)| = 10
\end{equation}
Innerhalb des Tubus $T(\Sigma_8) \subset \mathbb{O}$ mit transversalem Radius $R_0 = 10$ ist der Raum frei von interferierenden Nullstellen-Defekten ($\Sigma_\rho$).
\end{theorem}

\section{Epistemische Schutzwand (Claim-Wall)}

\begin{quote}
\small\textbf{Wichtige Abgrenzung:} Die hier formulierte Aharonov-Bohm-Analogie ist eine rein eichtechnische und analytisch-topologische Konstruktion auf der Spektralebene der $E_8$-Zeta-Funktion. Die Polstelle bei $s=8$ erzeugt kein physikalisches elektromagnetisches Feld, keine Kraftwirkungen (wie Lorentz-Kräfte) auf realweltliche Masseladungen und begründet keine reale physikalische Dynamik außerhalb des mathematischen Formalismus.
\end{quote}

\end{document}
